\chapter{Test-and-Fallback in Depth}\label{test-and-fallback-in-depth}

The test-and-fallback algorithms are the adaptive robustness mechanism of Chapter 4 and the
object of the theory in Chapter 9. This chapter gives their first empirical study: the
robustness envelope they achieve (§6.1), a counter-intuitive failure of the acceptance
threshold (§6.2), its recalibration and the resolution limit that recalibration exposes
(§6.3) --- the empirical face of the impossibility we prove in Chapter 9 --- and a benchmark of
the dynamic combiner that explains the \emph{commit-once} structure (§6.4). Throughout, the
\(\ell_1\) test is the empirical surrogate the original authors also fall back to (Chapter 3).

\section{The robustness envelope}\label{the-robustness-envelope}

Sweeping advice error \(\ell_1(p,q)\) on few-types instances (\(n=2000\), \(r=8\), prefix \(k=200\),
40 trials; \textbf{Figure 6.1}), FollowPrediction degrades \emph{linearly} from \(1.000\) at perfect
advice to \(0.453\) at \(\ell_1\!\approx\!1.1\) --- well \emph{below} the advice-free floor
(\(\approx0.99\)). TestAndMatch instead stays on the \textbf{upper envelope}: it captures the
benefit when advice is good and, when advice is bad, its prefix test rejects it and it falls
back to Ranking, never crashing (BEM \(0.998\to0.969\); Choo \(1.000\to0.991\)). This is the
adaptive counterpart of the structural robustness of Chapter 4.

\section{A threshold that is too lenient on average-case inputs}\label{a-threshold-that-is-too-lenient-on-average-case-inputs}

Sweeping the prefix (testing) size \(k\) at \emph{borderline} advice (\(\eta=0.15\), true
\(\ell_1\!\approx\!0.16\)) --- where the test works hardest (\textbf{Figure 6.2}) --- a larger, more
accurate test makes the \emph{worse} decision: as \(k\) grows \(25\to800\), the ratio \emph{falls}
\(0.992\to0.956\) and the misjudgement rate \emph{rises} \(0.00\to0.60\). The mechanism: the Choo/BEM
threshold \(\tau\) is calibrated to the worst-case baseline \(\beta\approx0.696\), but on these
instances the baseline is \(\approx0.99\) (F3), so the empirical break-even sits at
\(\ell_1\approx0\), far below \(\tau\). A small noisy prefix over-estimates \(\ell_1\) and
\emph{accidentally rejects} the borderline advice (landing safely on the floor); a large accurate
prefix correctly measures \(\ell_1\approx0.16<\tau\) and \emph{accepts} the mildly-bad advice the
worst-case threshold deems acceptable --- underperforming the baseline. On strong-baseline
inputs the worst-case threshold is too lenient, and a more accurate test only follows it more
faithfully.

\section{Recalibration, and the resolution limit it exposes}\label{recalibration-and-the-resolution-limit-it-exposes}

Recalibrating \(\tau\) to the \emph{measured} baseline \(\hat\beta\) eliminates the pathology: at
borderline advice the worst-case threshold's misjudgement climbs \(0.03\to1.00\) as the prefix
grows and its ratio drops to \(0.920\), while the recalibrated threshold holds misjudgement
\(0.00\) at every prefix and ratio \(\approx0.986\). But recalibration exposes a deeper limit.
At perfect advice the worst-case threshold scores \(1.000\) while the recalibrated one scores
only \(0.987\): the recalibrated \(\tau\approx2(1-\hat\beta)\approx0.028\) is \emph{smaller than the
empirical-\(\ell_1\) estimator's noise floor} (\(\approx0.05\)--\(0.13\)), so it can never
confidently \emph{accept} --- it rejects everything, including perfect advice, and always plays the
baseline. In short:

\begin{quote}
On strong-baseline instances, no practical empirical-\(\ell_1\) threshold can both capture
the consistency upside and stay safe: the upside is tiny and sits \emph{below the estimator's
resolution}. The worst-case threshold over-accepts; the recalibrated one over-rejects; a
better tester only follows whichever threshold more faithfully.
\end{quote}

This is exactly the phenomenon Chapter 9 proves unavoidable --- there, for \emph{any} test on a
sublinear prefix, not only the empirical-\(\ell_1\) threshold. This chapter is its empirical
shadow; Chapter 9 is the theorem.

\section{The dynamic combiner is dominated, and shows why matching needs test-then-commit}\label{the-dynamic-combiner-is-dominated-and-shows-why-matching-needs-test-then-commit}

We benchmark the Chłędowski-style dynamic combiner to contextualize the commit-once structure
of test-and-fallback. In its robust tuning the combiner sits exactly on the floor (\(0.990\)
across all advice quality): it never crashes but captures none of the consistency upside,
strictly dominated by TestAndMatch. More instructively, an \emph{eager} combiner that switches
mid-stream reveals a penalty specific to irrevocable problems: with perfect advice, eager
switching scores \(0.927\) --- \emph{below both} the pure follower (\(1.000\)) and the pure baseline
(\(0.958\); \texttt{tests/test\_combiner\_small.py}) --- because switching from Ranking to Mimic mid-run
lands the committed matching in an \emph{incompatible hybrid}. In an irrevocable problem the
follow/fallback decision must be made \emph{before} the bulk of the commitments, which is why
Choo/BEM test a prefix and then \emph{commit} rather than switching dynamically. The dynamic
combiner that is cheap insurance for caching does not port cleanly to matching.

\section{Chapter summary}\label{chapter-summary}

Test-and-fallback delivers the adaptive robustness envelope of §6.1, but its accept/reject
decision is fundamentally limited on strong-baseline inputs: no empirical-\(\ell_1\) threshold
both captures the upside and stays safe (§6.3), and dynamic switching is dominated by
test-then-commit (§6.4). The resolution limit of §6.3 is the empirical shadow of the
impossibility theorem of Chapter 9.
